% Appendix -- referenced from main body (Appendix A/B/C). Pages: not counted
% against the 8-page main limit.
\appendix

\section{Open-Source Release}
\label{app:repo}

The complete system source, this paper's \LaTeX{} source, and the seed
\texttt{paper\_meta/references.bib} that the lit-review sub-agent
verified are released under the MIT License at the public GitHub
repository below. Runtime manager/worker JSONL under \texttt{state/} is
local to each deployment (and \texttt{gitignore}d by default); integers
in \S\ref{sec:case} are reproduced from shell/git commands and verifier
artefacts named in Table~\ref{tab:engineering} (including its reproducibility
note), not from
a bundled export of those logs.

\begin{center}
  \url{https://github.com/guoshaoyang-pku/shaoyang-autoresearch}
\end{center}

\paragraph{Public mirror vs.\ source-branch history.}
The public GitHub checkout uses a squashed \texttt{main} history for
readability. Line-count and smoke-test integers in
Table~\ref{tab:engineering} are pinned to that released tree with the
shell commands printed under the table. Fine-grained commits under
\texttt{tools/cursor\_manager} on the authors' LiuLab branch (where this
paper was dogfooded), e.g.\ the dispatcher defect surfaced at
\texttt{c9ceeef}, are reproducible from that branch's git log but are not
individually replayable from the squashed public counter alone.

\noindent\textbf{Pinned SHAs for Table~\ref{tab:engineering}.}
\textbf{(i)} Public checkout: clone the repository above, checkout
\texttt{main}, and run \texttt{git rev-parse HEAD} for the exact commit
used in the \texttt{wc} and smoke-test rows (reported as
\texttt{POST\_RELEASE\_COMMIT\_OSS} in the camera-ready artefact list).
\textbf{(ii)} Source-branch aggregates: on the authors' LiuLab checkout,
on \texttt{feature/autoresearch-sub-agents}, run the \texttt{git log} invocations
printed beneath Table~\ref{tab:engineering} at \texttt{git rev-parse HEAD}
(meaning: the integers in the table are exact only for the branch tip that
contains this \texttt{paper\_meta} snapshot).

\noindent An internal corporate mirror may track the same files under a
different host; the GitHub URL above is the canonical citeable artefact.

\section{Sub-Agent Contract Template}
\label{app:contract-template}

The contract template that allowed four sub-agent modules to ship in
parallel (\S\ref{sec:arch-contract}, Table~\ref{tab:engineering})
follows the schema below. Each instance of the schema is a single
Markdown section in \texttt{docs/sub\_agents\_contract.md}; the four
sections together total \(593\) lines (\texttt{wc -l docs/sub\_agents\_contract.md} on the released tree).

\begin{quote}\small
\noindent\textbf{Module \texttt{<name>}.} \\
\textbf{Files I may create or modify:} \texttt{lib/<name>.py},
  \texttt{prompts/<name>.md}, \texttt{state/<name>/} (runtime-only). \\
\textbf{Files I must NOT touch:} all sibling modules' files (named
  explicitly), the parent \texttt{mgr} CLI, \texttt{lib/local\_worker.py}. \\
\textbf{Public Python API:} \texttt{class <Name>:} with methods
  \texttt{<sig>} returning dataclass \texttt{<Name>Report}; raises
  \texttt{<ExceptionList>}. \\
\textbf{Output schema:} writes
  \texttt{state/<name>/<wid>/<run\_id>.json} with keys
  \texttt{<keys>} and \texttt{state/<name>/<wid>/<run\_id>.md} with
  sections \texttt{<sections>}. \\
\textbf{Self-test:} \texttt{python -m lib.<name> --self-test} must
  exit 0 in \(<10\)~s with no network. \\
\textbf{Parent \texttt{mgr} wiring:} subcommand \texttt{<name>}
  with arguments \texttt{<args>}, exit codes \texttt{<map>}, stdout
  format \texttt{<format>}.
\end{quote}

\section{Lit-Review Verification Log}
\label{app:litreview-table}

\begin{table}[t]
  \caption{Outcome of \texttt{mgr lit-review paper\_c} on the seed
    bibliography of this paper. ``Defect found'' means the seed entry
    had a non-trivial metadata error; the dispatcher raised a
    \texttt{[BLOCKED-DECISION]}-equivalent advisory rather than
    silently rewriting the bibliography.}
  \label{tab:litreview}
  \centering
  \footnotesize
  \begin{tabular}{l c c p{0.38\textwidth}}
    \toprule
    \textbf{Citation key} & \textbf{Verified} & \textbf{Defect} & \textbf{Defect detail (if any)} \\
    \midrule
    \texttt{song2026paperorchestra}        & yes & no  & --- \\
    \texttt{sibyl2026system}               & yes & yes & seeded with fictitious ``renamed from FARS'' note; corrected to acknowledge separate provenance \\
    \texttt{fars2026analemma}              & yes & yes & added in second pass after the FARS-vs-us positioning gap was identified \\
    \texttt{yamada2025aiscientistv2}       & yes & no  & --- \\
    \texttt{karpathy2026autoresearch}      & yes & yes & canonical repo name is lowercase \texttt{autoresearch}, not \texttt{AutoResearch} \\
    \texttt{schmidgall2025agentlaboratory} & yes & no  & --- \\
    \texttt{lu2024aiscientist}             & yes & no  & --- \\
    \texttt{tang2025airesearcher}          & yes & no  & --- \\
    \texttt{miyai2026paperreconstruction}  & yes & no  & --- \\
    \bottomrule
  \end{tabular}
\end{table}

\section{Sentinel Incident Reconstruction and Lark Notification}
\label{app:sentinel-msg}

\paragraph{Incident timeline.}
On 2026-05-02 between approximately \texttt{14:35+0800} and
\texttt{14:50+0800}, the deployment host \texttt{gpu\_develop} began
producing identical exit-code-1 failures from the LLM CLI in use at
the time on every five-minute scheduler tick. (The CLI was
cursor-agent under an early prototype configuration; the same failure
mode is reproducible on the current codex-only loop, since the silent
escalation-log buildup is a property of the scheduler, not of the
underlying CLI.) The manager LLM dutifully wrote one
\texttt{tick\_cli\_failed} entry per tick to
\texttt{state/escalations.jsonl} but had no facility to surface the
pattern out of band. The streak continued unobserved for
approximately twenty-four hours; the deployment owner discovered it
when manually inspecting the escalation log on 2026-05-03.
Root-cause analysis identified a stale CLI session token; the
resolution was a single re-authentication followed by a manager
session reset (the codex-only loop equivalent is
\texttt{./kickoff.sh -{}-reset <wid>} after the underlying provider
proxy is healthy again).

\paragraph{Counterfactual under Sentinel.}
The Sentinel watchdog (\texttt{sentinel.py}, 580 lines) is
parameterised so that the same-cause-streak heuristic
(\S\ref{sec:discipline-sentinel}, item~ii) fires after
\textit{four} consecutive identical \texttt{exit\_code} values within
a configurable window. Under this parameterisation the same incident
would have triggered a Lark notification at approximately
\texttt{14:55+0800} -- twenty minutes after the first failure and
before the third hour of the streak -- yielding a wall-clock
detection-latency reduction by a factor of approximately
\(24 \text{h} / 20 \text{min} \approx 72\).

\paragraph{Rendered Lark message.}
The notification body is generated by \texttt{lib/lark\_notify.py}
from a hand-coded template. The dedupe key in the footer is the
SHA-1 prefix used by Sentinel to suppress duplicate alerts within a
four-hour window.

\begin{quote}\small\ttfamily
\textbf{[cursor\_manager:sentinel] same-cause failure streak detected} \\
worker: paper\_a \\
host: gpu\_develop \\
streak: 4 consecutive ticks with exit\_code=1 \\
last\_ticks: \\
\quad 2026-05-02T14:35:00+0800  exit=1  reason=tick\_cli\_failed \\
\quad 2026-05-02T14:40:00+0800  exit=1  reason=tick\_cli\_failed \\
\quad 2026-05-02T14:45:00+0800  exit=1  reason=tick\_cli\_failed \\
\quad 2026-05-02T14:50:00+0800  exit=1  reason=tick\_cli\_failed \\
suggested action: ssh gpu\_develop; codex --version; ./mgr backend-test;
tail -200 state/escalations.jsonl \\
\rule{\linewidth}{0.4pt} \\
dedupe\_key: tick\_cli\_failed:paper\_a:exit\_1 \\
suppressed\_until: 2026-05-02T18:50:00+0800
\end{quote}

\section{Full Limitations List}
\label{app:limitations}

\S\ref{sec:limitations} lists four scope limitations in main-body
form. Two further limitations were moved here for space.

\paragraph{(L5) Lit-review claim alignment.} The lit-review
sub-agent verifies that a cited work exists on Semantic Scholar but
does not verify that the cited work supports the surrounding claim
-- the same gap exists in \citet{song2026paperorchestra}. Closing
this gap would require a second LLM call per citation that reads
both the local paragraph and the cited work's abstract; we have not
budget-evaluated that addition.

\paragraph{(L6) Sentinel and Lark coupling.} Sentinel's four
trigger patterns are hand-tuned and frozen; a more principled
approach would learn them from a longer deployment history.
\texttt{lark\_notify.py} wraps an internal Lark CLI; teams outside
ByteDance must replace this adapter (the public API consists of one
function with one positional argument).

\section{Reproducibility Checklist}
\label{app:reproducibility}

We follow the NeurIPS 2026 reproducibility checklist. The system source
and this paper's source are released together (Appendix~\ref{app:repo})
under MIT license. Integer rows in \S\ref{sec:case} tie to explicit
commands in Table~\ref{tab:engineering} and its caption ---
\texttt{state/} JSONL is deployment-local unless an operator commits it
deliberately. The lit-review
sub-agent's Semantic Scholar queries used the public unauthenticated
endpoint and are reproducible up to API rate limits; the dispatcher
emits a structured advisory rather than a number when an entry cannot
be re-verified, so partial reruns are detectable. The cross-model
ablation (\S\ref{sec:case-ablation}) requires an API-reachable codex CLI
and provider credentials; we list the exact model strings and the proxy
configuration in \texttt{docs/codex\_backend\_recipe.md}. No
human-subject data is used; no model is trained; no GPU is occupied.
